% !TEX program = xelatex
\documentclass[11pt,a4paper]{article}

\usepackage[a4paper,top=2cm,bottom=2.2cm,left=2.2cm,right=2.2cm]{geometry}
\usepackage{fontspec}
\setmainfont{Times New Roman}
\setmonofont{Consolas}[Scale=0.92]

\usepackage{amsmath, amssymb}
\usepackage{booktabs}
\usepackage{xcolor}
\usepackage{tcolorbox}
\usepackage{enumitem}
\usepackage{titlesec}
\usepackage[hidelinks]{hyperref}

\definecolor{usthblue}{HTML}{0054A6}
\definecolor{usthdark}{HTML}{003366}
\definecolor{qbox}{HTML}{FFF8E6}
\definecolor{qborder}{HTML}{E5C97A}

\titleformat{\section}{\Large\bfseries\color{usthblue}}{Q\thesection.}{0.5em}{}
\titleformat{\subsection}{\normalsize\bfseries\color{usthdark}}{}{0pt}{}
\titlespacing*{\section}{0pt}{20pt}{8pt}
\titlespacing*{\subsection}{0pt}{12pt}{4pt}

\newtcolorbox{question}{
  colback=qbox, colframe=qborder, boxrule=0.6pt,
  arc=4pt, left=10pt, right=10pt, top=6pt, bottom=6pt,
  title={\textit{Question}}, fonttitle=\bfseries\color{black!70},
  before skip=8pt, after skip=8pt
}
\newtcolorbox{answer}{
  colback=blue!4, colframe=usthblue!50, boxrule=0.6pt,
  arc=4pt, left=10pt, right=10pt, top=6pt, bottom=6pt,
  title={\textit{Answer}}, fonttitle=\bfseries\color{usthblue},
  before skip=8pt, after skip=8pt
}
\newtcolorbox{tip}{
  colback=gray!5, colframe=gray!40, boxrule=0.4pt,
  arc=3pt, left=8pt, right=8pt, top=4pt, bottom=4pt,
  title={\footnotesize \textit{Delivery tip}}, fonttitle=\bfseries\color{black!60},
  before skip=6pt, after skip=6pt
}

\setlength{\parskip}{4pt}
\setlength{\parindent}{0pt}

\title{
  \vspace{-1.5cm}
  \color{usthblue}\Large\textbf{Defense Q\&A Preparation}\\[4pt]
  \color{black!75}\normalsize The ten hardest expected questions + model answers
}
\author{Nguyen Trong Bach --- 23BI14057}
\date{}

\begin{document}
\maketitle
\vspace{-0.8cm}

{\hypersetup{linkcolor=black}\tableofcontents}
\newpage

% =============================================================
\section{Why soft voting instead of stacking or learned weights?}

\begin{question}
``You chose soft voting to average four model probabilities. Why not learn ensemble weights (stacking, a meta-learner)? Would that not be better?''
\end{question}

\begin{answer}
I considered all three options and chose soft voting for \textbf{three practical reasons}:

\textbf{First}, stacking needs a second validation set to fit the meta-learner. If I reuse the same val set already used for checkpoint selection, the meta-learner overfits the \textit{quirks} of that val set rather than learning real signal. The ISIC 2018 split is already small; taking a slice for stacking means removing it from train, which is not worth the trade-off.

\textbf{Second}, hard voting (each model takes argmax then votes) discards probability information. Soft voting keeps the full vector $p_m(x) \in [0,1]^7$, which lets the ensemble distinguish ``all four models are very sure of MEL'' from ``two models vote MEL 51\%, two vote NV 49\%''. This distinction matters clinically.

\textbf{Third}, soft voting already reaches 89.49\% in-distribution and 90.49\% with metadata. McNemar's test shows soft voting beats every single model with $p < 10^{-7}$ against the strongest backbone. There is no clear \textit{headroom} for stacking to break through, and stacking adds inference cost.
\end{answer}

\begin{tip}
If the committee pushes back with ``I still think stacking would be better'' --- concede that it might, but say ``I prioritised reproducibility and simplicity at a result that already matches SOTA. Stacking is the next direction if more time and data become available.''
\end{tip}

% =============================================================
\section{How do you know your Grad-CAM is ``correct''?}

\begin{question}
``A red-and-yellow heatmap looks pretty to anyone. How do you know the model is genuinely looking at the lesion and not at some artefact that happens to overlap?''
\end{question}

\begin{answer}
I do not stop at qualitative inspection. I measure \textbf{four quantitative metrics} against a lesion mask (generated by Otsu thresholding) on the \textit{entire test set}:

\begin{itemize}[leftmargin=*]
  \item \texttt{focus\_ratio} --- fraction of total activation inside the lesion mask. High = correct focus. All four models reach $> 0.6$.
  \item \texttt{peak\_distance} --- distance from peak activation to lesion centroid (normalised by image diagonal). Low = focus on the centre. All four models reach $< 0.2$.
  \item \texttt{entropy} --- Shannon entropy of the heatmap. Low = sharp focus.
  \item \texttt{coverage\_75} --- \% area containing the top-25\% activation. Low = compact.
\end{itemize}

This is an improvement over many works that only show qualitative heatmaps without measurement. I learned the importance of quantification the hard way: my first attempt used \texttt{model.bn2} for EfficientNet and \texttt{focus\_ratio} was only 0.18 --- the heatmaps were uniformly blue. I switched to \texttt{blocks[-1][-1]} and \texttt{focus\_ratio} jumped to 0.51 --- only then did the heatmaps actually outline the lesion. Without the metric, the \texttt{bn2} mistake would have slipped past visual inspection.
\end{answer}

\begin{tip}
Emphasise the ``quantitative measurement caught a bug'' story --- it shows scientific thinking, not just running code.
\end{tip}

% =============================================================
\section{How do Macro F1 and balanced accuracy differ? Why choose Macro F1?}

\begin{question}
``You report Macro F1 rather than balanced accuracy. How do they differ? Why this choice?''
\end{question}

\begin{answer}
\textbf{Balanced accuracy} = mean of per-class recall. It looks only at recall and ignores precision.

\textbf{Macro F1} = mean of per-class F1 = mean of the harmonic mean of \textit{both} precision and recall.

\textbf{Practical difference:} balanced accuracy only penalises \textit{misses} (false negatives). Macro F1 also penalises \textit{false alarms} (false positives). On data with 58:1 imbalance, a model can achieve high balanced accuracy by predicting DF for every ambiguous image --- DF recall reaches 100\%, balanced accuracy rises, but DF precision collapses. Macro F1 catches that.

I chose Macro F1 because both error types have clinical cost: missing MEL is dangerous, but false alarms also trigger unnecessary biopsies, raising costs and stressing patients.

I still report Sensitivity ($=$ Recall) and Specificity ($=$ 1 $-$ false positive rate) per class for the complete picture.
\end{answer}

% =============================================================
\section{Why does Swin-Tiny overfit faster than the CNNs?}

\begin{question}
``Swin-Tiny early-stops at epoch 13 while the CNNs train to epochs 39--49. Why?''
\end{question}

\begin{answer}
Three cumulative reasons:

\textbf{(1) Lack of inductive bias.} CNNs come with translation invariance and locality built in --- convolution kernels enforce them by construction. A Transformer learns both from data, so it needs more data to converge. ISIC 2018 has only 7{,}010 training images --- small by computer-vision standards (ImageNet has 1.28M).

\textbf{(2) Extremely high parameter-to-sample ratio.} Swin-Tiny has 28M parameters and the training set has 7{,}010 images, giving a ratio of \textbf{4{,}000$\times$}. That is extremely overfit-prone. EB0's 5.3M parameters give 756$\times$, much more reasonable.

\textbf{(3) Quadratic self-attention.} Every pair of patches attends to every other pair, making it easy to memorise specific spatial structures of the training images.

I mitigated as much as I could: learning rate 1e-4 (CNNs use 3e-4) to preserve the pretrained attention weights, weight decay 5e-2 (CNNs use 1e-3) for stronger regularisation, and a 5-epoch warmup (CNNs use 3). Early stopping with patience 12 catches the epoch-13 checkpoint.

The silver lining: Swin's \textit{multi-seed mean} accuracy is the highest of any single model (88.33\% vs.\ EB0 87.31\%). Early stopping does not mean low performance --- it means \textit{higher variance}. In the ensemble, high variance combined with uncorrelated errors against the CNNs makes Swin a valuable ensemble component.
\end{answer}

% =============================================================
\section{What happens on images outside the 7 classes (e.g.\ acne)?}

\begin{question}
``You only trained on 7 classes. If a patient photographs acne, a scar, or some non-lesion --- what does the model do?''
\end{question}

\begin{answer}
The model is \textit{forced} to assign one of the 7 classes because the softmax outputs sum to 1 --- there is no ``other'' class. This is the genuine limitation of any closed-set classifier.

I integrated two mitigations:

\textbf{Confidence threshold.} The Streamlit app uses a 50\% threshold; if the top-1 probability is below 50\%, it shows a ``Low confidence'' warning and suggests consulting a specialist. In the ISIC test set the low-confidence rate is only $\sim$2\%, so the warning does not spam the user.

\textbf{Predictive entropy.} The app also computes the entropy of the probability vector. High entropy (near $\log 7$) means the model is unsure between many classes --- a possible OOD signal. The app displays entropy numerically.

A more complete production solution would be a dedicated \textit{open-set classifier} or \textit{OOD detector} (e.g.\ Mahalanobis distance on features, or an Energy-based score). I have not implemented this --- it is the obvious next step.
\end{answer}

\begin{tip}
This is a question a clinical or AI-focused examiner will care about. Acknowledge the limitation honestly and explain mitigation. Saying ``the model will classify into the closest class'' is wrong.
\end{tip}

% =============================================================
\section{How many false-negative MEL cases are there? Clinical impact?}

\begin{question}
``You report a melanoma F1 of 0.744. Concretely, out of 167 MEL test images, how many does the model miss? What is the clinical impact?''
\end{question}

\begin{answer}
\textbf{Numbers on the ISIC 2018 test set:}
\begin{itemize}[leftmargin=*]
  \item Test set contains 167 actual MEL images.
  \item Ensemble correctly identifies 124 $\to$ \textbf{43 false negatives} (recall 0.740).
  \item Among the 43 misses: 35 are misclassified as NV (benign nevus), 8 as BKL.
\end{itemize}

\textbf{Clinical impact if deployed as-is:} for every 100 actual MEL patients, 74 would be correctly flagged, 26 would be labelled ``benign'' and risk leaving the clinic without a biopsy. This is exactly why the Streamlit prototype \textbf{never} replaces a clinician --- it is a second opinion.

\textbf{Mitigation is available:} threshold tuning with Youden's J (slide A18) lowers the MEL decision threshold from 0.143 (default) to 0.100, raising MEL Sensitivity from 0.74 to 0.83 --- misses drop from 26/100 to 17/100. The trade-off is more false positives (specificity drops from 0.97 to 0.95 --- acceptable clinically, since missing MEL is much more dangerous than a false alarm).
\end{answer}

% =============================================================
\section{Why not Focal Loss instead of Label Smoothing?}

\begin{question}
``Focal Loss is designed for imbalanced data. Did you try it? Why pick Label Smoothing?''
\end{question}

\begin{answer}
I did try both. The \texttt{src/models/losses.py} file contains a \texttt{FocalLoss} class --- ready to use. My reasons for choosing Label Smoothing in the final pipeline:

\textbf{Focal Loss} ($\mathrm{FL}(p) = -(1-p)^\gamma \log p$) increases the weight on \textit{hard samples} within a batch. In a setup where I already have Weighted Sampler balancing classes plus Mixup/CutMix producing soft mixed labels, Focal Loss creates \textit{complex interactions}: a Mixup-mixed sample has a soft label like $[0, ..., 0.7, ..., 0.3]$, and Focal Loss does not have a canonical formula for soft labels --- I would have to hand-craft one. Result: F1 fluctuates and is less stable than cross-entropy plus label smoothing.

\textbf{Label Smoothing} is simpler: it multiplies the one-hot label by $(1-\varepsilon)$ and spreads $\varepsilon/C$ across all classes. PyTorch's standard cross-entropy supports soft labels natively (\texttt{label\_smoothing} argument), so it composes cleanly with Mixup/CutMix. The ablation shows removing Label Smoothing drops F1 by 3.11\%.

\textbf{In short:} the ``Weighted Sampler + Mixup/CutMix + Label Smoothing'' recipe is enough to handle imbalance (final 89.5\% Acc, 0.845 F1) without needing Focal Loss. If forced to choose just one of the two \textit{without} Mixup, Focal Loss might win.
\end{answer}

% =============================================================
\section{What is your ECE? Is the model calibrated?}

\begin{question}
``You report a high Macro AUC of 0.98. But AUC only measures ranking, not calibration. What is your ECE?''
\end{question}

\begin{answer}
I measured Expected Calibration Error (ECE) for all 4 models and applied temperature scaling (Guo et al., 2017):

\begin{center}\small
\begin{tabular}{lcc}
\toprule
\textbf{Model} & \textbf{ECE before} & \textbf{ECE after temperature scaling} \\
\midrule
EfficientNet-B0 & 0.041 & 0.018 \\
ResNet50        & 0.052 & 0.022 \\
DenseNet121     & 0.039 & 0.019 \\
Swin-Tiny       & 0.067 & 0.024 \\
\bottomrule
\end{tabular}
\end{center}

\textbf{Before calibration:} ECE 4--7\% --- the models are mildly overconfident, especially Swin (6.7\%). When the model says ``MEL 90\%'', the actual accuracy in that confidence bucket is only $\sim$83\%.

\textbf{After temperature scaling:} ECE drops to 2--2.4\% --- the probabilities are now trustworthy enough for clinical decisions. Ensembling calibrated models lowers ECE further.

Code: \texttt{src/utils/calibration.py} and \texttt{src/calibration\_analysis.py}.
\end{answer}

% =============================================================
\section{The dataset is mostly fair-skinned. Demographic bias impact?}

\begin{question}
``HAM10000 was collected in Austria and the USA --- predominantly fair skin (Fitzpatrick I--III). If you deploy in Vietnam (Fitzpatrick III--V), what happens?''
\end{question}

\begin{answer}
This is a genuine limitation of my thesis, which I acknowledge and discuss explicitly. Three points:

\textbf{(1) Skin-tone bias is a known field-wide issue.} Adamson \& Smith (2018) showed that commercial dermoscopy AI (such as SkinVision) under-performs significantly on darker skin --- MEL recall drops by $\sim$20\% on Fitzpatrick V--VI versus I--II. The field does not yet have a comprehensive solution.

\textbf{(2) I do not yet have Vietnamese data to test directly.} But the OOD evaluation on PAD-UFES-20 (Brazil, with more Fitzpatrick III--IV than HAM10000) shows the model \textbf{drops 56 pp in accuracy}. A meaningful portion of that drop \textit{may} come from skin-tone shift; another portion comes from modality shift (smartphone vs dermoscope). I cannot disentangle the two contributions.

\textbf{(3) Plausible remediation paths:}
\begin{itemize}[leftmargin=*]
  \item Fine-tune on Vietnamese data once available ($\sim$500 images would suffice).
  \item Skin-tone-aware colour augmentation: shift hue/saturation according to the Fitzpatrick scale.
  \item Group fairness loss (reweight the loss by skin-tone subgroup).
\end{itemize}

I do not claim the model is ready for Vietnam without adaptation. The Streamlit prototype is a \textit{research prototype} with a clear disclaimer.
\end{answer}

\begin{tip}
This is an \textbf{ethics question}. Answer by \textit{accepting responsibility}, not by arguing or deflecting blame.
\end{tip}

% =============================================================
\section{What more does production deployment require?}

\begin{question}
``If a hospital wanted to actually deploy this system, what else is needed beyond what you have built?''
\end{question}

\begin{answer}
I have thought through the steps:

\textbf{1. Regulatory compliance.} Medical diagnostic software needs CE marking (EU MDR) or FDA 510(k) clearance --- the prototype does not. Equivalent systems like SkinVision and MoleScope went through this process over 1--2 years.

\textbf{2. In-domain data.} As discussed in Q9, $\sim$500--5000 images from the actual hospital, taken by its own staff, are needed for fine-tuning to reduce domain shift.

\textbf{3. Open-set / OOD detection.} As discussed in Q5, the system needs a component that flags inputs that are \textit{not} skin lesions (acne, scars, wrong-region photos).

\textbf{4. Clinical UI/UX.} My Streamlit prototype is a demo. Production needs:
\begin{itemize}[leftmargin=*]
  \item PACS / EHR integration.
  \item An audit log of every prediction for traceability when errors occur.
  \item A clear workflow defining who has final responsibility (clinician).
\end{itemize}

\textbf{5. Post-deployment monitoring.} Data drift detection (does the input distribution change after six months?), performance monitoring (does accuracy degrade?), a feedback loop for clinicians to correct mistakes.

\textbf{6. Liability and insurance.} When the model errs and harms a patient, who is responsible? The clinician? The hospital? The vendor? Contracts must clarify this before deployment.

I do not claim this thesis is production-ready. It is a \textit{research prototype} demonstrating technical feasibility. Real deployment requires $\sim$2 years of engineering plus clinical trials.
\end{answer}

% =============================================================
\section{Defense kit checklist}

\textit{A list of items to prepare before defense day.}

\subsection{Hardware}
\begin{itemize}[leftmargin=*, itemsep=2pt]
  \item Primary laptop (battery full + charger).
  \item USB-C / HDMI adapter for the projector.
  \item External mouse or presentation clicker.
  \item Backup USB containing: all slides + thesis PDF + study guide PDF + demo images.
  \item Small headphones (in case you want to replay a rehearsal recording).
\end{itemize}

\subsection{Software \& demo}
\begin{itemize}[leftmargin=*, itemsep=2pt]
  \item Streamlit demo: run \texttt{streamlit run app.py} 15 minutes before defense to warm up.
  \item Open the browser tab at \texttt{http://localhost:8501}.
  \item Set theme (Dark/Light) via menu $\vdots$ $\to$ Settings $\to$ Theme.
  \item Upload one demo image and view Grad-CAM for all 4 models to cache them.
  \item Have 3 demo images ready: 1 clear MEL, 1 benign NV, 1 hard case (low-confidence to demonstrate the warning).
\end{itemize}

\subsection{Backup plans}
\begin{itemize}[leftmargin=*, itemsep=2pt]
  \item If Streamlit crashes: 5 demo screenshots (prediction + Grad-CAM) embedded in the slides.
  \item If the projector does not recognise the laptop: PDF slides on USB, open them on the room's machine.
  \item If internet fails: app.py runs 100\% offline, no internet needed.
\end{itemize}

\subsection{Printed materials}
\begin{itemize}[leftmargin=*, itemsep=2pt]
  \item Slides 4-up handout for the committee to follow.
  \item One printed copy of the English thesis.
  \item Printed study guide for your own reference (not required to hand to the committee).
\end{itemize}

\subsection{One hour before defense}
\begin{itemize}[leftmargin=*, itemsep=2pt]
  \item Test on the actual projector (if allowed).
  \item Listen to the last rehearsal recording.
  \item Light snack (banana, almonds) --- avoid hunger during the talk.
  \item Restroom.
  \item Phone to airplane mode.
\end{itemize}

\end{document}
